\documentclass[11pt]{article}
\usepackage[T1]{fontenc}
\usepackage{lmodern}
\usepackage[margin=1in]{geometry}
\usepackage{amsmath,amssymb,amsthm,mathtools}
\usepackage{array,booktabs,microtype}
\usepackage[hidelinks]{hyperref}
\numberwithin{equation}{section}
\newtheorem{theorem}{Theorem}[section]
\newtheorem{proposition}[theorem]{Proposition}
\newtheorem{lemma}[theorem]{Lemma}
\newtheorem{corollary}[theorem]{Corollary}
\theoremstyle{definition}
\newtheorem{definition}[theorem]{Definition}
\theoremstyle{remark}
\newtheorem{remark}[theorem]{Remark}
\newcommand{\HH}{\mathcal H}
\newcommand{\ceil}[1]{\left\lceil #1\right\rceil}
\newcommand{\floor}[1]{\left\lfloor #1\right\rfloor}
\newcommand{\pos}[1]{\left(#1\right)_+}
\newcommand{\bt}{\frac67}
\allowdisplaybreaks[2]
\setlength{\emergencystretch}{2em}
\title{A six-sevenths bound for transversals in hypergraphs with property $(7,2)$}
\author{}
\date{}
\begin{document}
\maketitle
\begin{abstract}
A hypergraph has property $(7,2)$ if every subfamily of at most seven
edges has a transversal of size at most two. We prove that every
$k$-uniform hypergraph with this property satisfies
$\tau(\HH)\le\lceil6k/7\rceil$ for $k\ge8$.
This improves the previously established asymptotic upper coefficient
$7/8$ to $6/7$. The proof combines sharper local avoidance constructions
with three successive exclusions of pair-intersection sizes and a
maximal-intersection argument. All constructions admit integral
allocations of points among at most four requests.
\end{abstract}

\section{Introduction}
A \emph{transversal} of a hypergraph is a set meeting every edge, and its
\emph{transversal number} $\tau(\HH)$ is the minimum size of such a set.
For integers $k\ge2$ and $r\ge3$, let $f(k,r)$ be the largest transversal
number of a finite $k$-uniform hypergraph in which every subfamily of at
most $r$ edges has a transversal of size at most two. The use of ``at most''
avoids vacuous exceptions for families with fewer than $r$ edges.

Erd\H{o}s, Fon-Der-Flaass, Kostochka and Tuza studied these local
transversal conditions in~\cite{EFKT}. For $r=7$, Fon-Der-Flaass,
Kostochka and Woodall proved the bounds
\[
 \ceil{3k/4}\le f(k,7)\le\ceil{7k/8}
\]
for the upper-bound range $k\ge8$; see~\cite{FKW}.
The question whether $f(k,7)=(3/4+o(1))k$ remains the motivating problem.
Here we improve the coefficient in the upper bound to $6/7$.

\begin{theorem}\label{thm:main}
Let $k\ge8$. Every finite $k$-uniform hypergraph with property $(7,2)$
satisfies
\[
 \tau(\HH)\le\ceil{6k/7}.
\]
Consequently $\limsup_{k\to\infty}f(k,7)/k\le6/7$.
\end{theorem}

Our proof refines the avoidance method of Fon-Der-Flaass, Kostochka
and Woodall~\cite{FKW}, which already uses good triples, intersection
gaps, and a maximal small intersection. The improvement comes from
stronger local completion constructions and their three-stage
combination. Broader local-cover problems were subsequently studied
by Kostochka~\cite{Kostochka} and Buci\'c, Kor\'andi and Sudakov~\cite{BKS}.

The proof uses the elementary consequence of $\tau(\HH)>T$ that each
set of at most $T$ points is avoided by an edge of $\HH$. Starting from
three edges with empty common intersection, we construct four such
requests so that no pair of points can meet all seven resulting edges.
Some constructions use a previously established gap in the possible
pair-intersection sizes. This permits three successive exclusions,
after which a pair with largest intersection at most $k/2$ supplies the
final contradiction.

The main argument is given in Sections~\ref{sec:gaps} and~\ref{sec:finish}.
The elementary setup is in Section~\ref{sec:prelim}. For completeness,
Appendices~\ref{app:allocation}--\ref{app:smallmax} contain full proofs of
all local constructions used in the argument. The local statements keep the integer request budget explicit.
The gap endpoints are chosen from the actual integer budget
$\lceil6k/7\rceil$, so no additive rounding term is needed.

\begin{table}[ht]
\centering\small
\begin{tabular}{>{\raggedright\arraybackslash}p{0.30\textwidth}>{\raggedright\arraybackslash}p{0.58\textwidth}}
\toprule
Part of the proof & Local results used\\
\midrule
First gap, Proposition~\ref{prop:gapone}
 & Lemmas~\ref{lem:asymone}, \ref{lem:fourcase}, \ref{lem:symmetric}\\
Middle gap, Proposition~\ref{prop:gaptwo}
 & Lemmas~\ref{lem:hallstatic} and~\ref{lem:gaptwo}\\
Final gap, Proposition~\ref{prop:gapthree}
 & Lemmas~\ref{lem:split}, \ref{lem:asymtwo},
   \ref{lem:gapempty}, \ref{lem:gapsurvive}\\
Local finisher, Proposition~\ref{prop:localfinish}
 & Lemmas~\ref{lem:asymone}, \ref{lem:fourcase}, \ref{lem:symmetric},
   \ref{lem:gaptwo}, \ref{lem:nearcore}\\
Small-maximum case, Lemma~\ref{lem:smallmax}
 & Lemmas~\ref{lem:balanced}, \ref{lem:smalltriple}, \ref{lem:twolarge}\\
\bottomrule
\end{tabular}
\caption{Dependencies of the main argument. The local results are proved
in full in the appendices; the balanced request is proved below.}
\label{tab:dependencies}
\end{table}

\section{Avoidance requests and good triples}\label{sec:prelim}
Throughout the main argument, $\HH$ is a finite $k$-uniform hypergraph
with property $(7,2)$ and $\tau(\HH)>T$, where $0\le T\le k$ is an integer. An \emph{avoidance request} is a set $D$ of at most $T$ points.
There is an edge disjoint from $D$, called a response to the request.
We construct responses successively. The requests need not be disjoint,
and their sizes are bounded individually.

A triple $E,F,G$ is \emph{good} if $E\cap F\cap G=\varnothing$.
Write
\[
 X=E\cap F,\qquad Y=E\cap G,\qquad Z=F\cap G,
\]
and write $P_E,P_F,P_G$ for the private parts of its three edges.
The six cells $X,Y,Z,P_E,P_F,P_G$ are disjoint. Any pair meeting the
three edges either lies in two distinct pair cells or lies in a pair
cell and the opposite private part. Points outside $E\cup F\cup G$
are irrelevant, because their partner alone cannot meet a good triple.
This observation verifies the static constructions below.

For four edges with all triple intersections empty, a pair meeting all
four must lie in two complementary pair cells. If precisely one triple
cell survives, there is in addition the product of that cell with the
opposite edge. These descriptions will be used explicitly in the local
proofs.

Call a subfamily \emph{bad} if its transversal number is greater than
two.

\begin{lemma}[Request-cover principle]\label{lem:requestcover}
Suppose $\tau(\HH)>T$. Let $A_1,\ldots,A_t$ be edges and let
$D_1,\ldots,D_s$ be sets of size at most $T$, where $t+s\le7$.
If every transversal of $\{A_1,\ldots,A_t\}$ of size at most two is
contained in some $D_j$, then responses avoiding the $D_j$ together
with the initial edges form a bad subfamily of at most seven edges.
\end{lemma}
\begin{proof}
A transversal of size at most two for all resulting edges would meet
the initial edges, and hence be contained in one of the requests. It
would then miss that request's response. Repeated responses only reduce
the number of distinct edges and do not affect the contradiction.
\end{proof}

We will also use integral allocation among requests. Suppose disjoint
sets of points may each be assigned to a specified subset of request
indices, and each request has an integer remaining capacity. An
assignment exists if every collection of source sets has total size at
most the sum of the capacities of its available indices. To see this,
make a bipartite flow network with the source-set sizes and request
capacities on its outer arcs. These inequalities are precisely its
min-cut conditions; integral augmenting paths give an assignment of
individual points. In particular, a set available to every request can
be distributed whenever its size is at most the sum of their remaining
capacities. For restricted assignments we list all Hall inequalities.

For integers $0\le L\le H<k$, let $\mathcal G(L,H)$ denote the condition
\[
 |A\cap B|\le H\quad\Longrightarrow\quad |A\cap B|\le L
 \qquad(A,B\in\HH, A\ne B).
\]
This formulation fixes both endpoints of a global intersection gap.
No ordering or convexity assumption on the set of possible intersections
is implicit in this notation.

\begin{lemma}[Balanced request]\label{lem:balanced}
Suppose $E,F$ have intersection size $M\le T\le k$. There is a request
of size $T$ containing $E\cap F$ such that an avoiding response $G$
forms a good triple and
\[
 |E\cap G|,|F\cap G|\le k-\floor{(T+M)/2}.
\]
\end{lemma}
\begin{proof}
Include the $M$ common points and distribute the remaining $T-M$
requested points as evenly as possible between $E\setminus F$ and
$F\setminus E$. Both selections fit because $T\le k$. Each new trace
has size at most $k-M-\lfloor(T-M)/2\rfloor$, which is the asserted
quantity.
\end{proof}

\begin{lemma}[Two capped traces]\label{lem:caps}
Suppose two edges meet in $q$ points. If a nonnegative integer $K$
satisfies
\[
 0\le2k-T-q\le2K,\qquad K\le k-q,
\]
there is a size-$T$ request whose response completes the pair to a
good triple with both new pair intersections at most $K$.
\end{lemma}
\begin{proof}
Choose integers $u,v\in[0,K]$ with $u+v=2k-T-q$. Retain $u$ and $v$
points in the two private parts, and request every other point of the
union of the edges. The retained sets fit because $K\le k-q$.
The request has size $2k-q-u-v=T$ and contains the common intersection.
Its response has the required traces and forms a good triple.
\end{proof}

All pair sizes used below are actual integer cardinalities.
A triple \emph{closes at budget $T$} if four further requests of size
at most $T$ force a bad subfamily of at most seven edges.
Each closing result is incompatible with property $(7,2)$ under
$\tau(\HH)>T$.

\section{Integer endpoints and three intersection gaps}\label{sec:gaps}
Fix $k\ge8$ and put
\begin{equation}\label{eq:budget}
 h=\floor{k/7},\qquad s=k-7h,\qquad T=k-h=\ceil{6k/7}.
\end{equation}
Thus $h\ge1$, $0\le s\le6$, and $6k/7\le T\le k$.
Assume $\tau(\HH)>T$. Define integer endpoints
\begin{equation}\label{eq:endpoints}
 \begin{gathered}
 A=\floor{T/2},\qquad B=\floor{(4T-2k)/3},\qquad
 C=T-\ceil{k/2},\\
 D=2k-T-2B,\qquad L=D-1,\qquad U=2T-k-\ceil{k/2}.
 \end{gathered}
\end{equation}
The first stage excludes intersections in $(A,B]$, up to $\lfloor k/2\rfloor$.
If this already reaches $k/2$, the finishing theorem applies. Otherwise
the second stage excludes $[D,C]$, and the third excludes $(B,k/2]$.
All endpoints and all retained subsets are integral.

\begin{lemma}[Endpoint identities]\label{lem:endpoints}
The endpoints in \eqref{eq:endpoints} satisfy
\begin{gather}
 0\le C\le A\le B,\qquad B\ge3h,\qquad
 L\le U=C-h,\qquad4L\le T,\label{eq:basicendpoints}\\
 k-B-1+L\le T.\label{eq:lastendpoint}
\end{gather}
If $B<\lfloor k/2\rfloor$, then $h<D\le C$; in particular,
$0\le L<C<k$.
\end{lemma}
\begin{proof}
Put $r=\lfloor(h+2s)/3\rfloor$. Directly from \eqref{eq:endpoints},
\begin{gather*}
 A=3h+\floor{s/2},\qquad B=3h+r,\qquad
 C=\floor{(5h+s)/2},\\
 D=2h+s-2r,\qquad L=2h+s-2r-1,\qquad
 U=\floor{(3h+s)/2}=C-h.
\end{gather*}
Clearly $C\ge0$, while $C\le A$ follows from $T\le k$.
Also $(4T-2k)/3\ge T/2$ because $T\ge4k/5$, so $B\ge A$.
The formula for $B$ gives $B\ge3h$.

To prove $L\le U$, it suffices to show
$2r+1\ge\lceil(h+s)/2\rceil$.
For every nonnegative integer $n$, writing $n=3a+b$ with $0\le b\le2$
gives $n\le4a+2$, and hence
$\lceil n/2\rceil\le2\lfloor n/3\rfloor+1$.
Apply this with $n=h+s$ and use $r\ge\lfloor(h+s)/3\rfloor$.
Next, write $h+2s=3r+t$ with $0\le t\le2$. Then
\[
 4L-T=2t-2r-s-4\le0.
\]
Finally,
\[
 k-B-1+L-T=s-3r-2\le-h-s\le0,
\]
because $3r\ge h+2s-2$. This proves
\eqref{eq:basicendpoints}--\eqref{eq:lastendpoint}.
If $B<\lfloor k/2\rfloor$, then $2B<k$ implies $D>k-T=h$.
Also $D=L+1\le C-h+1\le C$, since $h\ge1$.
\end{proof}

The cap requests used in the three stages are summarized in
Table~\ref{tab:gaps}. When a stage is reached, its earlier exclusions
apply to every pair of family edges, including pairs involving a new
response.
\begin{table}[ht]
\centering
\begin{tabular}{cccc}
\toprule
Stage & Hypothetical intersection $q$ & Retained cap & New traces\\
\midrule
1 & $A+1\le q\le\min(B,\lfloor k/2\rfloor)$ & $C$ & $\le C$\\
2 & $D\le q\le C$ & $B$ & $\le A$\\
3 & $B+1\le q\le\lfloor k/2\rfloor$ & $C$ & $\le L$\\
\bottomrule
\end{tabular}
\caption{Integer gap exclusions. Stages 2 and 3 are needed only when
$B<\lfloor k/2\rfloor$.}\label{tab:gaps}
\end{table}

\subsection{The first gap}
\begin{proposition}\label{prop:gapone}
No pair intersection has integer size
$A+1\le q\le\min(B,\lfloor k/2\rfloor)$.
\end{proposition}
\begin{proof}
Use Lemma~\ref{lem:caps} with retained cap $C$.
Its host bound follows from $C\le\lfloor k/2\rfloor\le k-q$, and
$2k-T-q\ge0$. Furthermore
\[
 2k-T-2C=3h+(k\bmod2)\le A+1\le q,
\]
so the request has size $T$. Its good triple has pair sizes
$x=q$ and $0\le z\le y\le C$. Since $C\le A<x$, the largest is $x$.
Put $d=y-z$, $v=y+z$, and $S=x+v$.

\paragraph{Case 1: $d>x-h$.}
Apply Lemma~\ref{lem:asymone} in orientation $(x,y,z)$.
As $v+d=2y\le2T-k$, its four conditions follow from
\begin{align*}
 S&<T,&k/2+y&\le T,\\
 (k+2x-d)/2&<(k+x+h)/2\le T,\\
 (k+2x+2v-d)/3&<(2k+T-x)/3\le T.
\end{align*}
For the third bound use
$x\le B\le(4T-2k)/3\le3T-2k$, where the last comparison follows from
$T\ge4k/5$. For the fourth, $x\ge A+1\ge3h+1>2h$.

\paragraph{Case 2: $d\le x-h$ and $v\le3T-k-2x$.}
Use Lemma~\ref{lem:fourcase} in orientation $(y,z,x)$.
Its nine expressions are bounded as follows:
\begin{gather*}
 v\le2T-k\le T,\qquad k/2+y, k/2+z\le T,\\
 k+d-x\le k-h=T,\qquad k-d-x\le k-x\le T,\\
 k-S/3\le T,\qquad (3k+S)/5\le T,\\
 (k+v+2x)/3\le T,\qquad (2k+3x)/4\le T.
\end{gather*}
Here $S\ge x\ge3h$ proves the bound on $k-S/3$.
For the next expression use $S\le x+2T-k$ and $x\le3T-2k$.
The last two use the case assumption and $x\le B$.

\paragraph{Case 3: $v>3T-k-2x$.}
Use Lemma~\ref{lem:symmetric}. Its four exact budget forms satisfy
\[
 \begin{gathered}
 k+x-v<2k+3x-3T\le T,\qquad (k+x)/2\le T,\\
 (2k+x+d)/3<T,\qquad(3k+S)/5\le T.
 \end{gathered}
\]
The second follows from $x\le B\le2T-k$. For the third,
$d\le2T-k-v$ gives $2k+x+d<2k-T+3x\le3T$.
The fourth is the estimate in Case 2. Every triple therefore closes.
\end{proof}

If $B\ge\lfloor k/2\rfloor$, Proposition~\ref{prop:gapone} already
excludes every intersection greater than $A$ and at most $k/2$.
Theorem~\ref{thm:finish} then gives the contradiction.
For the remaining two stages assume
\begin{equation}\label{eq:branch}
 B<\floor{k/2}.
\end{equation}

\subsection{The second gap}
\begin{proposition}\label{prop:gaptwo}
Under \eqref{eq:branch}, no pair intersection has size $D\le q\le C$.
\end{proposition}
\begin{proof}
The interval is nonempty and positive by Lemma~\ref{lem:endpoints}.
Use retained cap $B$ in Lemma~\ref{lem:caps}.
The hosts fit because $q\le C\le\lfloor k/2\rfloor$ and
$B<\lfloor k/2\rfloor$. The required total is nonnegative and at most
$2B$ because $q\ge D=2k-T-2B$.
By the first gap, the two new traces are at most $A$.
Thus there is a good triple with
\[
 D\le x\le C,\qquad0\le z\le y\le A.
\]

If $x+z\ge k-B$, apply Lemma~\ref{lem:gaptwo} in orientation $(y,x,z)$
with $\mathcal G(A,B)$, supplied by the first gap.
The host conditions follow from $x+y\ge x+z\ge k-B$.
The first request has size
\[
 \max(x,2k-x-2B)\le T,
\]
since $x\ge D$. The remaining conditions are
\begin{gather*}
 y+A, z+A\le2A\le T,\qquad k+2x\le k+2C\le2T,\\
 3k\le4T,\qquad3k+S\le3k+C+2A\le5k/2+2T\le5T.
\end{gather*}

If $x+z<k-B$, use Lemma~\ref{lem:hallstatic} with distinguished pair
$a=y$ and other pairs $b=x,c=z$. We have $a\le A$, $b\le C$, and
\begin{align*}
 c&\le k-B-x-1\le k-B-D-1=T+B-k-1\le C,\\
 b+c&\le k-B-1\le3T-2k,
\end{align*}
where the last bound uses $B\ge3h$.
The four base capacities are nonnegative because
$a+b,a+c\le A+C\le T$ and $b+c\le k-B-1\le T$.
The seven Hall conditions follow from
\begin{gather*}
 k+2b, k+2c\le2T,\qquad k+3a\le k+3T/2\le3T,\\
 2k+b+c\le3T,\qquad
 2k+2a+b, 2k+2a+c\le3k/2+2T\le4T,\\
 3k+a\le3k+T/2\le4T.
\end{gather*}
Thus the exact integral allocation closes this case as well.
\end{proof}

\subsection{The third gap}
\begin{proposition}\label{prop:gapthree}
Under \eqref{eq:branch}, no pair intersection has size
$B+1\le q\le\lfloor k/2\rfloor$.
\end{proposition}
\begin{proof}
Use retained cap $C$. The same host calculation as in the first stage
applies, since $q\ge B+1\ge3h+1$.
By the second gap, both new traces are at most $L=D-1$. Thus the good
triple satisfies
\[
 B+1\le x\le\floor{k/2},\qquad0\le z\le y\le L.
\]
The global implication is $\mathcal G(L,C)$. Define
\[
 J=2k-C+2\floor{k/2}-3T=\floor{(h-s)/2},\qquad S=x+y+z.
\]

\paragraph{Case 1: $S\le2T-k$.}
Use Lemma~\ref{lem:split}, whose bounds follow from
\[
 (k+S)/2\le T,\qquad k-x+|y-z|\le k-B-1+L\le T,
 \qquad k/3+x\le5k/6\le T.
\]

\paragraph{Case 2: $2T-k\le S\le T$ and $z<J$.}
Use Lemma~\ref{lem:asymtwo} in orientation $(x,z,y)$.
The first three forms fit by $S\le T$, $k/3+x\le T$, and
$k-x+y\le T$, using \eqref{eq:lastendpoint}.
For the last form, $z<J\le h/2$, $y\le L\le2h+s$, and $2x\le k$ give
\[
 k+2x+3z+y\le2k+3h/2+2h+s=35h/2+3s\le18h+3s=3T.
\]
This case is empty if $J\le0$.

\paragraph{Case 3: $2T-k\le S\le T$ and $z\ge J$.}
Apply Lemma~\ref{lem:gapempty} with thresholds $(L,C)$.
Expanding its first condition gives
\[
 S+p+t=\max(S,k-C+y,k-C+z,2k-2C-x)\le T.
\]
The middle terms fit because $L\le U=C-h$, while the last fits because
$2k-2C-T=3h+(k\bmod2)\le B+1\le x$.
For its split condition,
\[
 \begin{aligned}
 k+3x+p+t=\max\{&k+3x, 2k-C+2x-y, 2k-C+2x-z,\\
                &3k-2C+x-y-z\}\le3T.
 \end{aligned}
\]
The first term is at most $5k/2\le3T$.
The next two fit by $y\ge z\ge J$ and $2x\le2\lfloor k/2\rfloor$.
For the last,
\begin{align*}
 3k-2C+x-y-z
 &=3k-2C+2x-S\\
 &\le4k-2C+2\floor{k/2}-2T=6k-4T\le3T.
\end{align*}
The equality uses $C=T-\lceil k/2\rceil$.
The third request fits because $y+z+2L\le4L\le T$.

\paragraph{Case 4: $S\ge T$.}
Apply Lemma~\ref{lem:gapsurvive} with thresholds $(L,C)$ and
$Q=S-T$. Its conditions follow from $L\le U=C-h$ and $4L\le T$:
\begin{align*}
 x+y&\le\floor{k/2}+L\le\floor{k/2}+C\le T,\\
 k-x-y&\le k-T+z\le h+L\le C,\\
 k-x-z+Q&=h+y\le C,\\
 2x+2Q+k-y-z
 &=4x+y+z+k-2T\\
 &\le4\floor{k/2}+2U+k-2T=2T-3(k\bmod2)\le2T,\\
 y+z+2L&\le4L\le T,\\
 2k+x-y+L+Q
 &=2k+2x+z+L-T\\
 &\le2k+2\floor{k/2}+2U-T=3T-2(k\bmod2)\le3T.
\end{align*}
This exhausts the cases.
\end{proof}

Together, the first and third gaps exclude every intersection in
$(A,k/2]$. The next section proves that this suffices for the contradiction.

\section{The maximal-intersection argument}\label{sec:finish}
The three gap exclusions reduce the proof to the following statement.
\begin{theorem}[Finishing theorem]\label{thm:finish}
Let $k\ge8$ and $T=\lceil6k/7\rceil$. Suppose that every pair of
 distinct edges in a finite $k$-uniform hypergraph with property $(7,2)$
has intersection at most $T/2$ or greater than $k/2$.
Then $\tau(\HH)\le T$.
\end{theorem}
We first give the local completion needed when the largest intersection
at most $k/2$ is greater than $2k/7$.

\begin{proposition}\label{prop:localfinish}
Let $6k/7\le T\le k$ be an integer, and let $M\le T/2$ be an integer.
Suppose every pair intersection in the family is at most $M$ or greater
than $k/2$. Every good triple whose pair sizes are
$M\ge y\ge z\ge0$ closes at budget $T$.
\end{proposition}
\begin{proof}
Put $h=k-T$, $c=T-k/2$, and $S=M+y+z$.
Here $0\le h\le c$. We use four cases. All numerical comparisons
below follow from $6k/7\le T\le k$.

\paragraph{Case 1: $y\le c$.}
Write $d=y-z$. If $d>M-h$, apply Lemma~\ref{lem:asymone} in orientation
$(M,y,z)$. Its first three budget forms satisfy
\[
 S<2y+h\le T,\qquad k/2+y\le T,\qquad
 (k+2M-d)/2<(k+M+h)/2\le k-T/4\le T.
\]
For the fourth,
\[
 \frac{k+2M+2(y+z)-d}{3}
 <\frac{k-M+4y+3h}{3}
 \le\frac{k}{3}+y+h\le\frac{5k}{6}\le T,
\]
Here the first weak inequality uses $M\ge y$, and the next uses $y\le c$.

Suppose next that $d\le M-h$ but $S<3h$.
Use Lemma~\ref{lem:asymone} in orientation $(z,y,M)$.
The first two forms fit because $S<3h\le T$ and $k/2+y\le T$.
The last two are at most $(k+3h)/2$ and $(k+9h)/3$, respectively,
and both are at most $T$.

In the remaining subcase, $d\le M-h$ and $S\ge3h$.
Use Lemma~\ref{lem:fourcase} in orientation $(z,y,M)$.
The inequality $d\le M-h$ implies $M\ge h$, so
\[
 k-d-M\le k-M\le T,\qquad k+d-M\le T.
\]
Also $y+z\le2c=2T-k\le T$, $k/2+y,k/2+z\le T$, and
$k-S/3\le T$. The remaining three forms satisfy
\begin{gather*}
 S\le M+2y\le5T/2-k,
 \qquad \frac{3k+S}{5}\le\frac{2k}{5}+\frac T2\le T,\\
 \frac{k+y+z+2M}{3}\le T,
 \qquad \frac{2k+3M}{4}\le\frac k2+\frac{3T}{8}\le T.
\end{gather*}
This proves Case 1.

\paragraph{Case 2: $y>c$ and $h<z\le c$.}
Apply Lemma~\ref{lem:gaptwo} in orientation $(M,z,y)$, with
thresholds $(M,\lfloor k/2\rfloor)$.
Its pair-cell cuts have size $g=(k-z-\lfloor k/2\rfloor)_+$.
They fit because $M+z\ge y+z>c+h=k/2$, so both host sums are at
least $\lceil k/2\rceil$.
The initial request fits since
\[
 z+2g=\max\{z,k-z+(k\bmod2)\}\le T;
\]
here the integer $z$ is at least $h+1$.
The remaining hypotheses follow from
\[
 2M\le T,\qquad M+y\le2M\le T,\qquad
 k+2z\le2T,\qquad3k\le4T,
\]
and
\[
 3k+S\le3k+2M+c\le5k/2+2T\le5T.
\]

\paragraph{Case 3: $y>c$ and $z>c$.}
Use Lemma~\ref{lem:symmetric}. Its four budget forms are at most,
respectively,
\[
 2k-3T/2,\qquad k/2+T/4,\qquad5k/6,\qquad
 (3k+3T/2)/5.
\]
Each is at most $T$. For the third bound use
$M+y-z\le T-c=k/2$; the others follow directly from $M\le T/2$
and $y,z>c$.

\paragraph{Case 4: $y>c$ and $z\le h$.}
Use Lemma~\ref{lem:nearcore} in orientation $(M,z,y)$, with global
small-intersection cap $m=M$. Its common hypotheses follow from
\begin{gather*}
 y+z\le T/2+h=k-T/2\le T,\qquad M+y\le2M\le T,\\
 k+z-y<k+h-c=5k/2-2T\le T,\\
 2k+M-2y<3k-3T/2\le2T,\\
 2k+2M-y<5k/2\le3T,\qquad2M\le T.
\end{gather*}
If $S\le T$, use the empty-core variant. Its two further conditions are
\[
 k-M+z\le k-y+z<T,\qquad
 2k-2M+y\le2k-M<2k-c=5k/2-T\le2T.
\]
If $S>T$, put $\Delta=S-T$. After substitution, the four remaining
near-core conditions reduce to
\begin{align*}
 k+y+2z&\le3k-3T/2\le2T,\\
 2k-M+2y+z&\le2k+y+z\le3k-T/2\le3T,\\
 2k+M+z&\le3k-T/2\le3T,\\
 3k+y&\le3k+T/2\le4T.
\end{align*}
Thus both variants close the triple. Since $h\le c$, the four cases
exhaust all possibilities.
\end{proof}

\begin{proof}[Proof of Theorem~\ref{thm:finish}]
Suppose $\tau(\HH)>T$. Choose an edge avoiding a $T$-subset of another
edge. Their intersection has size at most $k-T\le k/2$, so the largest
pair-intersection size $M\le k/2$ exists. By hypothesis, $M\le T/2$.
If $M\le2k/7$, Lemma~\ref{lem:smallmax} applies.
Its four arithmetic conditions hold for the present budget.
For $k\ge14$, the estimate $T\ge6k/7$ gives respective margins at least
\[
 k/7-1,\qquad k/14-1/2,\qquad3k/14-1/2,\qquad2k/7-4,
\]
with the required strictness in the first. For $8\le k\le13$, we have
$T=k-1$, and the four conditions reduce to
$k>7$, $k-3\ge\lceil k/2\rceil$, $k\ge5$, and $k\ge8$.

Otherwise, choose a pair attaining $M>2k/7$ and apply
Lemma~\ref{lem:balanced}. The resulting good triple has other
intersections $Y\ge Z$ satisfying
\[
 Y,Z\le k-\floor{(T+M)/2}<3k/7+1/2\le k/2.
\]
By maximality, $Y,Z\le M$. Every pair intersection in the family is
at most $M$ or greater than $k/2$, so Proposition~\ref{prop:localfinish}
closes this triple. Both alternatives contradict property $(7,2)$.
\end{proof}

\begin{proof}[Proof of Theorem~\ref{thm:main}]
Suppose $\tau(\HH)>T=\lceil6k/7\rceil$.
If $B\ge\lfloor k/2\rfloor$, Proposition~\ref{prop:gapone} excludes
all intersections in $(A,k/2]$.
Otherwise, Propositions~\ref{prop:gapone}--\ref{prop:gapthree} do so.
Since $A=\lfloor T/2\rfloor$, every pair intersection is at most
$T/2$ or greater than $k/2$. Theorem~\ref{thm:finish} gives a
contradiction.
\end{proof}

\begin{corollary}
The conclusion of Theorem~\ref{thm:main} also holds for finite families
of nonempty sets of rank at most $k$.
\end{corollary}
\begin{proof}
Pad each edge to size $k$ with points private to that edge. Old
transversals still meet the padded edges. Conversely, replace each
private point of a transversal by an arbitrary point of its original
nonempty edge. This gives a transversal of the original family of no
larger size. Thus padding preserves the transversal number, and it
preserves property $(7,2)$ because every original two-point transversal
still works. Apply Theorem~\ref{thm:main} to the padded family.
\end{proof}

\appendix
\section{Integral allocations and a near-core construction}\label{app:allocation}
The request-cover and integral-allocation principles of
Section~\ref{sec:prelim} apply throughout this appendix.

\begin{lemma}[Exact boxed-triangle allocation]\label{lem:triangle}
Let $M_1,M_2,M_3$ be nonnegative integers and let $L_1,L_2,L_3$ be
integers satisfying $L_i\le M_j+M_l$ whenever $\{i,j,l\}=\{1,2,3\}$.
The minimum of $a_1+a_2+a_3$ over integral vectors satisfying
\[
 0\le a_i\le M_i,\qquad a_j+a_l\ge L_i
\]
is the ceiling of
\begin{equation}\label{eq:triangleexact}
 \max\left\{0,L_1,L_2,L_3,
 L_1+L_2-M_3,L_1+L_3-M_2,L_2+L_3-M_1,
 \frac{L_1+L_2+L_3}{2}\right\}.
\end{equation}
\end{lemma}
\begin{proof}
For a proposed nonnegative integer total $t$, the conditions are equivalent to
\[
 0\le a_i\le\min(M_i,t-L_i),\qquad \sum_i a_i=t.
\]
Integral coordinates exist if and only if all three upper bounds are
nonnegative and their sum is at least $t$. Expand the sum of the three
minima as the minimum of its eight affine expressions. The resulting
conditions are
\begin{gather*}
 t\ge0,L_1,L_2,L_3,\qquad
 t\ge L_i+L_j-M_l\quad(\{i,j,l\}=\{1,2,3\}),\\
 2t\ge L_1+L_2+L_3,\qquad t\le M_1+M_2+M_3,\qquad
 L_i\le M_j+M_l.
\end{gather*}
The last three host inequalities hold by assumption. Every lower bound
in \eqref{eq:triangleexact} is at most $M_1+M_2+M_3$, so the ceiling of
the maximum also satisfies the upper bound on $t$. Coordinates of this
total can be filled successively within their integral upper bounds.
The same argument without ceilings gives the real minimum.
\end{proof}

\begin{lemma}[Exact symmetric static allocation]\label{lem:symmetric}
Let a good triple have pair sizes $m\ge y\ge z\ge0$, and put $S=m+y+z$.
The four-request construction below can be implemented at integer
budget $T$ if and only if
\begin{equation}\label{eq:symexact}
 T\ge\max\left\{k+m-y-z,\frac{k+m}{2},
             \frac{2k+m+y-z}{3},\frac{3k+S}{5}\right\}.
\end{equation}
Whenever it can be implemented, it closes the triple.
\end{lemma}
\begin{proof}
Split the pair cells $M,Y,Z$ into selected parts of sizes $a,b,c$ and
their complements. With the private cells denoted by $P_E,P_F,P_G$,
use four requests
\begin{align*}
 R_0&=M_1\cup Y_1\cup Z_1,\\
 R_M&=M\cup P_G\cup Y_2\cup Z_2,\\
 R_Y&=Y\cup P_F\cup M_2\cup Z_2,\\
 R_Z&=Z\cup P_E\cup M_2\cup Y_2.
\end{align*}
Their sizes are $a+b+c$, $k+m-b-c$, $k+y-a-c$, and $k+z-a-b$.
Apply Lemma~\ref{lem:triangle} with hosts $(m,y,z)$ and requirements
\[
 L_1=k+m-T,\qquad L_2=k+y-T,\qquad L_3=k+z-T.
\]
The host conditions reduce to $T\ge k+m-y-z$. The individual
requirements $L_i\le T$ reduce to $2T\ge k+m$; the two-$L$ requirements
reduce to $3T\ge2k+m+y-z$; and the total requirement is
$5T\ge3k+S$. This proves the exact criterion.

Every pair meeting the good triple is contained in a request. A pair
cell and its opposite private part lie in the corresponding request.
For two different pair cells, if both points lie in selected parts they
belong to $R_0$; otherwise one of the last three requests contains both.
Thus four responses close the triple.

The optimality assertion concerns this specified construction only,
not all possible closing constructions.
\end{proof}

\begin{lemma}[Exact static Hall allocation]\label{lem:hallstatic}
For a good triple with pair sizes $a,b,c$, the four-request allocation
below has an integral implementation at budget $T$ if and only if
\begin{gather}
 T-a\ge0,\quad T-a-b\ge0,\quad T-a-c\ge0,\quad T-b-c\ge0,
 \label{eq:hallbase}\\
 k+2c\le2T,\quad k+2b\le2T,\quad k+3a\le3T,\quad
 2k+b+c\le3T,\label{eq:hallseven}\\
 2k+2a+c\le4T,\quad2k+2a+b\le4T,\quad3k+a\le4T.\nonumber
\end{gather}
Whenever it has such an implementation, it closes the triple.
\end{lemma}
\begin{proof}
Label the pair cells $X=E\cap F$, $Y=E\cap G$, $Z=F\cap G$,
of sizes $a,b,c$, by request-index sets $\{1,2,3\}$, $\{2,4\}$,
and $\{3,4\}$, respectively. A label specifies all requests containing
its point. Assign the private points to singleton labels according to
\[
 P_E:\{3\}\text{ or }\{4\},\qquad
 P_F:\{2\}\text{ or }\{4\},\qquad
 P_G:\{1\},\{2\}\text{ or }\{3\}.
\]
The four remaining capacities are exactly those in
\eqref{eq:hallbase}. The seven Hall inequalities for the three source
sets $P_E,P_F,P_G$, of sizes $k-a-b,k-a-c,k-b-c$, are precisely
\eqref{eq:hallseven}. Thus the integral-allocation principle proves the
claimed criterion.

Every candidate pair has intersecting labels. The three pair-cell
labels intersect pairwise, and every allowed label of a private cell
meets the label of its opposite pair cell. Consequently every candidate
pair is contained in some request. The request-cover principle gives
closure.

\end{proof}


\begin{lemma}[Near-core allocation]\label{lem:nearcore}
Suppose every pair intersection in the family is at most $m$ or greater
than $k/2$, where $m\le k/2$. Let $E,F,G$ be a good triple with pair
sizes $x=|E\cap F|$, $y=|E\cap G|$, $z=|F\cap G|$, and put
$\Delta=(x+y+z-T)_+$. The triple closes at integer budget $T$ if
$y+z\le T$ and
\begin{align}
 m+z&\le T,& k+y-z&\le T,\label{eq:ncfirst}\\
 2k+x-2z&\le2T,&2k+x+m-z&\le3T,\nonumber\\
 x+m&\le T,&&\nonumber\\
 k-x+y+\Delta&\le T,&2k-2x+z+\Delta&\le2T,\label{eq:nclast}\\
 2k-z+\Delta&\le2T,&3k-x-y+\Delta&\le3T.\nonumber
\end{align}
If $x+y+z\le T$, the last four bounds may instead be replaced by
\begin{equation}\label{eq:emptycore}
 k-x+y\le T,\qquad 2k-2x+z\le2T.
\end{equation}
\end{lemma}
\begin{proof}
Put $X=E\cap F$, $Y=E\cap G$, $Z=F\cap G$. Request an edge $H$
avoiding $Y\cup Z$ and a subset of $X$ of size
$\min(x,T-y-z)$. The only possibly nonempty triple cell is
$P=E\cap F\cap H$, of size $p\le\Delta$.
Define the disjoint cells
\begin{align*}
 X'&=X\setminus P, & A&=H\cap\bigl(E\setminus(F\cup G)\bigr),\\
 B&=H\cap\bigl(F\setminus(E\cup G)\bigr), &C&=H\cap G,\\
 W&=G\setminus(E\cup F\cup H).
\end{align*}
Write $a,b,c,w$ for the latter four sizes. Then
\begin{gather*}
 a\le k-x-y,\quad b\le k-x-z,\quad c\le k-y-z,\\
 p+a+b+c\le k,\qquad c+w=k-y-z.
\end{gather*}
The complete candidate products for pairs meeting $E,F,G,H$ are
\[
 P\times(Y\cup Z\cup C\cup W),\qquad
 X'\times C,\qquad Y\times B,\qquad Z\times A.
\]
A point in three edges lies in $P$; otherwise a covering pair must lie
in complementary pair cells. This proves completeness of the list.

\paragraph{Case 1: $p+a\le m$.}
Put $P$ in all three remaining requests, $X'$ in requests $1,3$,
$Y,B$ in request $1$, and $Z,A$ in request $2$. Partition $C$ between
requests $1,3$, and then partition $W$ among all three. The initial
loads are
\[
 L_1=x+y+b,\qquad L_2=p+z+a,\qquad L_3=x.
\]
They fit by $L_1\le k+y-z$, $L_2\le m+z$, and $L_3\le x+m$.
The first partition fits because
\[
 L_1+L_3+c=2x+y+b+c\le2k+x-2z\le2T.
\]
After that partition, the total load including $W$ is
\[
 L_1+L_2+L_3+c+w=2x+k+p+a+b\le2k+x+m-z\le3T.
\]
All capacities are integral, so the partitions exist. Each candidate
product is covered: $P$ lies in every request; the two possible labels
of $C$ meet the label of $X'$; and $Y,B$ and $Z,A$ share requests.

\paragraph{Case 2: $p+a>m$.}
The global dichotomy gives $p+a>k/2$. The traces $E\cap H$ and
$G\cap H$ are disjoint, so $c<k/2$ and hence $c\le m$.
Put $P$ in requests $1,2$, $X',C$ in request $2$, $Y,B$ in request
$1$, and $Z$ in requests $1,3$. Partition $A$ between requests $1,3$
and $W$ between requests $1,2$. The initial loads are
\[
 L_1=p+y+z+b,\qquad L_2=x+c,\qquad L_3=z.
\]
They fit by $L_1\le\Delta+k-x+y$, $L_2\le x+m$, and $L_3\le m+z$.
The three Hall conditions for the two flexible cells are
\[
 a\le2T-L_1-L_3,\qquad w\le2T-L_1-L_2,\qquad
 a+w\le3T-L_1-L_2-L_3.
\]
They follow from
\begin{align*}
 L_1+L_3+a&\le\Delta+2k-2x+z\le2T,\\
 L_1+L_2+w&=p+x+k+b\le\Delta+2k-z\le2T,\\
 L_1+L_2+L_3+a+w&=p+x+k+z+a+b
                  \le\Delta+3k-x-y\le3T.
\end{align*}
Thus an integral allocation exists. The labels of $P$ meet all allowed
labels of $Y,Z,C,W$; $X',C$ share request $2$; $Y,B$ share request $1$;
and both labels of $A$ meet the label $\{1,3\}$ of $Z$. Every candidate
pair is therefore missed by a response.

Finally, if $x+y+z\le T$, the first request omits the whole pair-cell
union, so $P=\varnothing$. The cell $W$ has no candidate-pair neighbour
and can be left out of every remaining request. In Case 2 only the
allocation of $A$ between requests $1,3$ is needed. Its conditions are
exactly the base-capacity bounds and \eqref{eq:emptycore}. Case 1 already
works even if $W$ is included. This proves the empty-core variant.
\end{proof}

% Complete local avoidance constructions. All mathematical statements
% are independent of the source-note numbering.
\section{Local avoidance constructions}\label{app:local}
Throughout this appendix, the family is $k$-uniform and has transversal
number greater than the integer request budget $0\le T\le k$. Thus every set
of at most $T$ points is avoided by a family edge. All edges obtained by
such requests belong to the original family. Repetitions cause no problem:
the contradiction is a subfamily of at most seven edges.

The notation of a good triple is as in Section~\ref{sec:prelim}.
A construction \emph{closes} a triple if four further requests force a
subfamily with no transversal of size at most two. Pair-cell sizes in
this appendix are actual integer cardinalities.

For integers $0\le L\le H<k$, the notation $\mathcal G(L,H)$ recalls
the global condition from Section~\ref{sec:prelim}: every intersection
of distinct family edges of size at most $H$ has size at most $L$.
Since $H<k$, any response with trace at most $H$ is distinct from the
edge against which that trace is measured.
The three threshold lemmas state their integer conditions directly.

\begin{lemma}[Three small intersections; cf.~\cite{FKW}]\label{lem:smalltriple}
Suppose all three pairwise intersections of a good triple have size at most $m\le k/2$. It extends to a bad subfamily of at most seven edges whenever
\[
\tau>B:=\left\lceil\max\{(3k+m)/4,(2k+2m)/3\}\right\rceil.
\]
\end{lemma}
\begin{proof}
Write the triple as $E,F,G$, with $X=E\cap G$, $Y=E\cap F$, $Z=F\cap G$, and relabel so $x=|X|\ge y=|Y|\ge z=|Z|$. Choose $P_G\subseteq G\setminus(E\cup F)$ of size $B-x-y$, and $B_0\subseteq F\setminus(E\cup G)$ of size $\min(B-x,k-y-z)$. Put $C=F\setminus(Y\cup B_0)$; then
\[
c=|C|=\max(k-B+x-y,z).
\]
Choose $P_E\subseteq E\setminus(F\cup G)$ of size $\min(B-x-c,k-x-y)$. All requested sizes are nonnegative and fit their hosts: $B-x-y\ge B-2m\ge0$, $c\le k-B+m$, $B-x-c\ge2B-k-2m\ge0$, and $B-x-y\le k-x-z$ because $y\ge z$ and $B\le k$.

Request edges avoiding
\[
X\cup Y\cup P_G,\quad X\cup B_0,\quad X\cup C\cup P_E,\quad (E\cup G)\setminus(X\cup P_E\cup P_G).
\]
The first three sizes are at most $B$. The fourth size is
\[
\max\{3(k-B)+x,\ 2(k-B)+y+z,\ (k-B)+2y\}\le B,
\]
by the definition of $B$.
A candidate pair with one point in $X$ has its other point in $Y,B_0$, or $C$ and is missed by one of the first three responses. Otherwise a point in $Y$ must pair with a point in $G\setminus X$, and the first or fourth response misses the pair according as the latter point is in $P_G$ or its complement. If the first point is in $Z$, use the third or fourth response and $P_E$. This exhausts the candidate pairs.
\end{proof}

\begin{lemma}[Splitting one pair cell]\label{lem:split}
For a good triple with pair sizes $x=|E\cap F|$, $y=|E\cap G|$, $z=|F\cap G|$, four legal responses yield a bad seven-tuple at any budget $T\le k$ satisfying
\[
T\ge\max\{(k+x+y+z)/2,\ k-x+|y-z|,\ k/3+x\}.
\]
The lemma is exactly integral.
\end{lemma}
\begin{proof}
Put $S=x+y+z$; the first hypothesis and $T\le k$ give $S\le T$. Request $H$ avoiding $X=E\cap F$, $Y=E\cap G$, $Z=F\cap G$, and
\[
p=\max(0,k-2(T-x)-y-z)
\]
points of the private part of $G$. This selection fits because $T\ge x$. Its union size is
\[
S+p=x+\max(y+z,k-2T+2x)\le T.
\]
All four triple cells vanish. Put $A=F\cap H$, $B=G\cap H$, $C=E\cap H$ with sizes $a,b,c$. Then
\[
a\le k-x-z,\quad b\le2(T-x),\quad c\le k-x-y,\quad a+b+c\le k.
\]

If $b\le T-x$, the three requests $X\cup B$, $Y\cup A$, $Z\cup C$ have sizes at most $T$: for the latter two use $k-x+|y-z|\le T$. They cover the three complementary candidate pairs.

If $b>T-x$, divide $B$ into two parts of size at most $T-x$; this is possible because $b\le2(T-x)$, also integrally. Two requests contain $X$ and one part each. The third contains $Y\cup Z\cup A\cup C$, whose size satisfies
\[
y+z+a+c\le k-b+y+z<k-T+x+y+z\le T.
\]
These requests again cover all three candidate pairs.
\end{proof}

\begin{lemma}[One surviving triple cell]\label{lem:fourcase}
Let $E,F,G$ be a good triple of $k$-sets, and put $X=E\cap F$, $Y=E\cap G$, $Z=F\cap G$, with sizes $x,y,z$ and $S=x+y+z$. Four further avoidance responses force a bad seven-tuple at any integer budget $T\le k$ satisfying
\[
\begin{gathered}
T\ge\max\Bigl\{x+y,\ k/2+x,\ k/2+y,\ k+x-y-z,\ k-x+y-z,\\
 k-S/3,\ (3k+S)/5,\ (k+x+y+2z)/3,\ (2k+3z)/4\Bigr\}.
\end{gathered}
\]
The construction is exactly integral. Any of the three pair cells can be designated $Z$.
\end{lemma}
\begin{proof}
Request $H$ avoiding $X,Y$ and a subset $Z_0\subseteq Z$ of size $\min(z,T-x-y)$. This is legal. Write
\[
Q=Z\cap H,\quad A=(F\cap H)\setminus Q,\quad B=(G\cap H)\setminus Q,\quad C=E\cap H,
\]
with sizes $q,a,b,c$. Also put $V=Z\setminus Q$, $v=z-q$, and $D=E\setminus(X\cup Y\cup C)$, $d=k-x-y-c$. These eight cells are disjoint, and
\[
q\le(S-T)_+,\quad a\le k-x-z,\quad b\le k-y-z,\quad c\le k-x-y,\quad q+a+b+c\le k.
\]
The only triple cell of $E,F,G,H$ is $Q$. A pair piercing these four edges therefore either belongs to $Q\times E$, or to one of $X\times B$, $Y\times A$, $V\times C$. We will make all three remaining avoidance requests contain $Q$ and their union contain $E$, while covering these last three products.

Two uniform bounds used below are
\[
q+x+b\le\max(k+x-y-z,k+2x-T)\le T,
\]
\[
q+y+a\le\max(k-x+y-z,k+2y-T)\le T.
\]
Also $z\le T$, since $4T\ge2k+3z$ and $z\le k$.

\paragraph{Case 0: $c\le T-z$.} Start the three requests with the disjoint-cell unions
\[
Q\cup X\cup B,\qquad Q\cup Y\cup A,\qquad Z\cup C.
\]
Each base has size at most $T$. Their total load, plus $d$, is
\[
k+2q+a+b+z\le\max(3k-S,3k+S-2T)\le3T.
\]
The first inequality follows from $a+b\le2k-x-y-2z$ and $q\le(S-T)_+$. Distribute $D$ among the remaining capacities. These capacities and $d$ are integers; hence the distribution can be integral. All candidate pairs are now covered by the bases or by $Q$ together with this distribution.

In the other cases $c>T-z$. Set
\[
a_0=k+x+z-2T,\qquad b_0=k+y+z-2T.
\]

\textbf{Case 1: $a<a_0$.} Use bases
\[
Q\cup X\cup B,\qquad Y\cup A\cup Z\cup C_2,\qquad Z\cup C_3,
\]
where $C=C_2\sqcup C_3$. This split fits because
\[
y+a+z\le k+x+y+2z-2T\le T,
\]
\[
a+c\le2k+z-y-2T\le2T-2z-y.
\]
Thus the two capacities $T-y-a-z$ and $T-z$ are nonnegative and sum to at least $c$. The total base load plus $d$ is
\[
k+q+a+b+2z\le2k+2z-c<2k+3z-T\le3T.
\]
Distribute $D$ among the remaining capacities. The three products are covered, and all bases contain $Q$.

\textbf{Case 2: $b<b_0$.} Interchange $X,B,x,b$ with $Y,A,y,a$ in Case 1. This covers the case even if both inequalities hold.

\textbf{Case 3: $a\ge a_0$ and $b\ge b_0$.} Put $u=c+z-T>0$. Since $z\le T$, we have $u\le c$, so choose $C_{12}\subseteq C$ of size $u$ and put $C_3=C\setminus C_{12}$. Choose an integer $t$ in
\[
\max(0,y+a+z+u-T)\ \le t\le\ \min(v,T-q-x-b-u).
\]
This interval is nonempty. The second upper entry is nonnegative because
\[
q+b+c\le k-a\le2T-x-z.
\]
The second lower entry is at most $v$ because
\[
q+a+c\le k-b\le2T-y-z.
\]
Finally the untruncated lower entry does not exceed the untruncated upper entry: this is equivalent to
\[
q+a+b+2c+x+y+3z\le4T,
\]
which follows from $q+a+b+c\le k$, $c\le k-x-y$ and $2k+3z\le4T$. All endpoints are integers.

Split $V=V_{13}\sqcup V_{23}$ with $|V_{13}|=t$, and use bases
\[
Q\cup X\cup B\cup C_{12}\cup V_{13},
\]
\[
Q\cup Y\cup A\cup C_{12}\cup V_{23},
\]
\[
Z\cup C_3.
\]
The interval inequalities make the first two sizes at most $T$; the third has size exactly $T$. Their total load plus $d$ is
\[
k+q+a+b+2z+u\le2k+3z-T\le3T.
\]
Again distribute $D$ into the residual capacities. The first two requests cover $X\times B$, $Y\times A$ and $V\times C_{12}$; the third covers $V\times C_3$. All contain $Q$ and their union contains $E$. Thus each candidate piercing pair of the first four edges misses at least one of the three new responses. This is a bad seven-tuple.
\end{proof}

\begin{lemma}[An asymmetric split]\label{lem:asymone}
With the good-triple notation of Lemma~\ref{lem:fourcase}, four further avoidance responses force a bad seven-tuple whenever $T\le k$ and
\[
T\ge\max\left\{S,\ \frac k2+y,\ \frac{k+2x-y+z}{2},\ \frac{k+2x+y+3z}{3}\right\}.
\]
The construction is exactly integral. All six permutations of $x,y,z$ are available.
\end{lemma}
\begin{proof}
Request $H$ avoiding $X,Y,Z$ and $T-S$ points in the private part of $F$. The private selection fits because $T-S\ge0$ and $T-S\le k-x-z$, the latter following from $T\le k+y$. Put $A=F\cap H$, $B=G\cap H$, $C=E\cap H$, with sizes $a,b,c$. All four triple cells vanish, so the candidate piercing pairs are $X\times B$, $Y\times A$, and $Z\times C$. Also
\[
a\le k-T+y,\qquad b\le k-y-z,\qquad a+b+c\le k.
\]
If $a\le T-y-z$, split $B=B_1\sqcup B_2$ into two parts of size at most $T-x-z$. This is possible, integrally, since
\[
2(T-x-z)\ge k-y-z\ge b.
\]
Start the final three requests with
\[
X\cup Z\cup B_1,\qquad X\cup Z\cup B_2,\qquad Y\cup Z\cup A.
\]
The first two fit by construction; the third fits by the case assumption. The total base load plus $c$ is
\[
2x+y+3z+a+b+c\le k+2x+y+3z\le3T.
\]
Distribute $C$ integrally among the residual capacities. Every final request contains $Z$, so this distribution covers $Z\times C$; the first two requests cover $X\times B$ and the third covers $Y\times A$.

If instead $a>T-y-z$, then
\[
b+c\le k-a<k-T+y+z\le2(T-x-z).
\]
Split $B\cup C$ into two parts of size at most $T-x-z$, and put $X\cup Z$ together with one part in each of the first two requests. The third request is $Y\cup A$, of size at most $k-T+2y\le T$. These three requests cover all three candidate-pair products again. Thus no pair pierces the seven edges in either case.
\end{proof}

\begin{lemma}[A second asymmetric split]\label{lem:asymtwo}
With the same good-triple notation, four further avoidance responses force a bad seven-tuple if $T\le k$ and
\[
T\ge\max\left\{S,\ \frac k3+x,\ k-x+z,\ \frac{k+2x+3y+z}{3}\right\}.
\]
The construction is exactly integral, with all six permutations available.
\end{lemma}
\begin{proof}
Request $H$ avoiding $X,Y,Z$ and
\[
p=\max(0,k-2(T-x)-y-z)
\]
points of the private part of $G$. This fits its host since $T\ge x$, and its total cost is $\max(S,k+3x-2T)\le T$. As in Lemma~\ref{lem:split}, write $A=F\cap H$, $B=G\cap H$, $C=E\cap H$, with sizes $a,b,c$. Then
\[
b\le2(T-x),\qquad c\le k-x-y,\qquad a+b+c\le k.
\]
All four triple cells vanish and the candidate pairs are $X\times B$, $Y\times A$, $Z\times C$.

If $b>k+y+z-T$, split $B$ into two parts of size at most $T-x$, and use $X$ with each part as two requests. The third is $Y\cup Z\cup A\cup C$, of size at most $k-b+y+z<T$.

Otherwise $b\le k+y+z-T\le2(T-x-y)$, where the last inequality is exactly $3T\ge k+2x+3y+z$. Split $B=B_1\sqcup B_2$ into two parts of size at most $T-x-y$. Start the three requests with
\[
X\cup Y\cup B_1,\qquad X\cup Y\cup B_2,\qquad Y\cup Z\cup C.
\]
The first two fit by construction; the third has size at most $k-x+z\le T$. The total base load plus $a$ is at most
\[
2x+3y+z+a+b+c\le k+2x+3y+z\le3T.
\]
Distribute $A$ into the residual integer capacities. Since every request contains $Y$, this covers $Y\times A$. The first two cover $X\times B$ and the third covers $Z\times C$. Both cases therefore yield a bad seven-tuple.
\end{proof}

\begin{lemma}[Two traces forced below a threshold]\label{lem:gapempty}
Suppose $\mathcal G(L,H)$ holds, and a good triple has pair sizes $x,y,z$.
Put $S=x+y+z$ and
\[
 p=(k-x-y-H)_+,\qquad t=(k-x-z-H)_+.
\]
The triple closes at integer budget $T\le k$ if
\begin{equation}\label{eq:g0exact}
 S+p+t\le T,\qquad k+3x+p+t\le3T,\qquad y+z+2L\le T.
\end{equation}
\end{lemma}
\begin{proof}
Choose $p$ private points of $E$ and $t$ private points of $F$ and
request them together with $X,Y,Z$. The private selections fit because
$H\ge0$. The initial size is $C_0=S+p+t\le T$. Pad the request by
$T-C_0$ private points of $G$; this fits because $C_0\ge S$ and $T\le k$.
Let $K$ be a response. Its traces on $E$ and $F$ are at most $H$, hence
at most $L$. All triple intersections among $E,F,G,K$ are empty.

Write $A=F\cap K$, $B=G\cap K$, $C=E\cap K$. Their sizes satisfy
$|A|,|C|\le L$ and $|B|\le k-T+x+p+t$.
Split $B$ into two nearly equal parts $B_1,B_2$. The requests
\[
 X\cup B_1,\qquad X\cup B_2,\qquad Y\cup Z\cup A\cup C
\]
have sizes at most
\[
 x+\ceil{(k-T+x+p+t)/2}\le T,\qquad y+z+2L\le T.
\]
The first bound is equivalent to $k+3x+p+t\le3T$, since all quantities
are integers. These requests cover the complementary products
$X\times B$, $Y\times A$, and $Z\times C$.
\end{proof}

\begin{lemma}[One surviving triple cell and a threshold]\label{lem:gapsurvive}
Suppose $\mathcal G(L,H)$ holds. For a good triple put $S=x+y+z$ and
$Q=(S-T)_+$. The triple closes at integer budget $T\le k$ if
\begin{align}
 x+y&\le T,& k-x-y&\le H,& k-x-z+Q&\le H,\label{eq:g1exact}\\
 2x+2Q+k-y-z&\le2T,& y+z+2L&\le T,&
 2k+x-y+L+Q&\le3T.\nonumber
\end{align}
\end{lemma}
\begin{proof}
Request $K$ avoiding $X,Y$ and a subset of $Z$ of size
$\min(z,T-x-y)$. Put
\[
 P=Z\cap K,\quad A=(F\cap K)\setminus P,\quad
 B=(G\cap K)\setminus P,\quad C=E\cap K,
\]
with sizes $p,a,b,c$. Then $p\le Q$, $b\le k-y-z$, and
\[
 c\le k-x-y\le H,\qquad a+p\le k-x-z+Q\le H.
\]
These are actual pair intersections. Hence $\mathcal G(L,H)$ gives
$c\le L$ and $a+p\le L$.

Split $B=B_1\sqcup B_2$ into nearly equal parts and start the final
three requests with
\[
 X\cup P\cup B_1,\qquad X\cup P\cup B_2,\qquad Y\cup Z\cup A\cup C.
\]
The first two sizes are at most
$x+Q+\lceil(k-y-z)/2\rceil\le T$ by the fourth inequality in
\eqref{eq:g1exact}. The third is at most $y+z+2L-p\le T$.
All bases contain $P$. To make their union contain $E$, distribute
$D=E\setminus(X\cup Y\cup C)$ among the remaining capacities.
The total load including $D$ is
\begin{align*}
 k+x+z+2p+a+b
 &\le k+x+z+p+L+b\\
 &\le2k+x-y+L+Q\le3T.
\end{align*}
An integral distribution therefore exists. The candidate products are
$P\times E$, $X\times B$, $Y\times A$, and $(Z\setminus P)\times C$.
The first is covered because every request contains $P$ and their union
contains $E$; the other three are covered by the stated bases.
\end{proof}

\begin{lemma}[Two large opposite pair cells]\label{lem:twolarge}
Suppose all pair intersections in the family are at most $m$ or greater than $k/2$, where $m\le k/4$. Let $E,F,G,H$ be four $k$-edges with all triple intersections empty. Put $X=E\cap F$, $B=G\cap H$, of sizes $x,b>k/2$, and suppose the other four pair cells each have size at most $m$. Three further legal responses give a bad seven-tuple if $T\le k$ and
\[
T\ge\lceil k/2\rceil+2m,\qquad T\ge x+m,\qquad 3T\ge2x+b+\lceil k/2\rceil+2m.
\]
\end{lemma}
\begin{proof}
Put $Y=E\cap G$, $Z=F\cap G$, $A=F\cap H$, $C=E\cap H$, and let $U=Y\cup Z\cup A\cup C$. These six pair cells are disjoint. Write $s=|Y|+|Z|$, $t=|A|+|C|$, so $s,t\le2m$, and define
\[
p=\lceil k/2\rceil-\min(s,t),\qquad q=T-\lceil k/2\rceil-\max(s,t).
\]
Both are nonnegative. Moreover $p\le\lceil k/2\rceil\le b$ and $q\le T-\lceil k/2\rceil\le\lfloor k/2\rfloor<x$. Choose $B_0\subseteq B$, $X_0\subseteq X$ with sizes $p,q$, and request $I$ avoiding $U\cup B_0\cup X_0$, whose size is exactly $T$.

The $G$-trace of $I$ has size at most
\[
k-s-p=\lfloor k/2\rfloor+\min(s,t)-s\le\lfloor k/2\rfloor,
\]
and the analogous statement holds for $H$. The global gap therefore forces both traces to have size at most $m$. In particular $d=|B\cap I|\le m$. Also $x_I=|X\cap I|\le x-q$.

Choose $B_1\subseteq B$ containing $B\cap I$ and having size $\min(b,T-x)$. This is possible since $T-x\ge m\ge d$. Use the last two requests
\[
X\cup B_1,\qquad (X\cap I)\cup(B\setminus B_1).
\]
The first has size at most $T$. For the second, note that
\[
x+x_I+b\le2x+b-q\le2x+b-T+\lceil k/2\rceil+2m\le2T.
\]
It follows that its size $x_I+\max(0,b-T+x)$ is at most $T$ as well.

A pair piercing $E,F,G,H$ belongs to $X\times B$, $Y\times A$ or $Z\times C$. The latter two products lie entirely in $U$ and hence miss $I$. For a remaining pair in $X\times B$, if its $B$-point belongs to $B_1$, the penultimate response misses both points. Otherwise its $B$-point lies outside $I$, since $B\cap I\subseteq B_1$; its $X$-point must then belong to $I$, and the last response misses both. This proves the lemma.
\end{proof}

\begin{lemma}[Two surviving triple cells]\label{lem:gaptwo}
Suppose $\mathcal G(L,H)$ holds. For a good triple put $S=x+y+z$ and
$g=(k-y-H)_+$. The triple closes at integer budget $T\le k$ if
\begin{gather}
 g\le\min(x,z),\qquad y+2g\le T,\label{eq:g2exact}\\
 x+L\le T,\quad z+L\le T,\quad k+2y\le2T,\quad
 3k\le4T,\quad 3k+S\le5T.\nonumber
\end{gather}
\end{lemma}
\begin{proof}
Choose $g$ points from each of $X,Z$ and request them together with $Y$.
Complete the part of the request inside $F$ to exactly $T-y$ points,
first using remaining points of $X\cup Z$ and then private points of
$F$. This is possible since $2g\le T-y\le k$. Let $K$ be a response.

Put $P=E\cap F\cap K$ and $Q=F\cap G\cap K$, of sizes $p,q$.
These are the only possible triple cells among the four edges because
$Y\cap K=\varnothing$. Set
\[
 A=(F\cap K)\setminus(P\cup Q),\quad
 B=(G\cap K)\setminus Q,\quad C=(E\cap K)\setminus P,
\]
with sizes $a,b,c$. Priority in the first request gives
$p+q\le(S-T)_+$, and its size inside $F$ gives
$d:=p+q+a=|F\cap K|\le k-T+y$.
The two $g$-point cuts ensure $|E\cap K|,|G\cap K|\le H$, so both
traces are at most $L$.

Start the last three requests with
\[
 X\cup(G\cap K),\qquad Z\cup(E\cap K),\qquad Y\cup(F\cap K).
\]
Their sizes are at most $x+L$, $z+L$, and $k-T+2y$, hence at most $T$.
Each contains both $P$ and $Q$.

The still-unassigned parts of $E\cup G$ are
$D=E\setminus(F\cup G\cup K)$ and
$W=G\setminus(E\cup F\cup K)$. They are disjoint and have sizes
$k-x-y-c$ and $k-y-z-b$. The total base load including these parts is
\[
 2k-y+2(p+q)+a=2k-y+(p+q)+d\le3k-T+(S-T)_+\le3T.
\]
The last inequality follows from $3k\le4T$ when $S\le T$ and from
$3k+S\le5T$ otherwise. Distribute $D\cup W$ integrally among the
remaining capacities, so the union of the requests contains $E\cup G$.

A candidate pair involving $P$ has its other point in $G$, while a pair
involving $Q$ has its other point in $E$. Such pairs are covered because
every request contains $P,Q$ and their union contains $E,G$. Every other
candidate pair lies in $(X\setminus P)\times B$,
$(Z\setminus Q)\times C$, or $Y\times A$, covered by the first, second,
or third base respectively.
\end{proof}


\section{Finishing when the largest small intersection is small}\label{app:smallmax}
\begin{lemma}\label{lem:smallmax}
Let $T$ be an integer with $6k/7\le T\le k$, satisfying
\begin{equation}\label{eq:fourinteger}
 \begin{gathered}
 6T>5k+1,\qquad 3T\ge2k+\ceil{k/2},\\
 2T\ge3k/2+1/2,\qquad 12T\ge10k+4.
 \end{gathered}
\end{equation}
If every pair intersection in a $k$-uniform hypergraph with property
$(7,2)$ is at most $2k/7$ or greater than $k/2$, then $\tau(\HH)\le T$.
\end{lemma}
\begin{proof}
Suppose $\tau(\HH)>T$. Choosing an edge avoiding a $T$-subset of
another gives an intersection at most $k-T$. Thus the largest pair
intersection $m\le k/2$ exists, and $m\le2k/7$. Choose a pair attaining
$m$ and use Lemma~\ref{lem:balanced}. Both new intersections are at most
$k-\lfloor(T+m)/2\rfloor$.

If both are at most $k/2$, then all three are at most $m$ by maximality.
Lemma~\ref{lem:smalltriple} applies at a budget at most $T$, since
\[
 (2k+2m)/3\le6k/7,\qquad (3k+m)/4\le23k/28<6k/7.
\]

Otherwise, relabel the good triple so that
\[
 x=|E\cap F|>k/2,\qquad |E\cap G|=m,\qquad |F\cap G|\le m.
\]
At most one of the new intersections is greater than $k/2$, since they
are disjoint subsets of the response edge. The balanced bound gives
\begin{equation}\label{eq:smallmaxbounds}
 x\le k-(T+m)/2+1/2,\qquad m\le k-T.
\end{equation}
For the second assertion, $T+m\ge k+1$ would imply
$\lfloor(T+m)/2\rfloor\ge\lceil k/2\rceil$, contrary to $x>k/2$.
Now $T\ge6k/7$ gives $m\le k/7<k/4$, and
\[
 x+2m\le(5k-4T+1)/2<T,
\]
since $6T>5k+1$. Request $H$ avoiding the three pair cells and enough
private points of $G$ to make the request size $T$. The padding fits:
the initial request size is $x+|E\cap G|+|F\cap G|$, and $T\le k$.
All triple intersections among $E,F,G,H$ are now empty. The $E$- and
$F$-traces of $H$ are less than $k/2$, because $x>k/2$, and are
therefore at most $m$. Write
\[
 b=|G\cap H|\le k-T+x.
\]
If $b\le k/2$, apply Lemma~\ref{lem:smalltriple} to $E,G,H$,
discarding $F$. The resulting bad subfamily has at most seven edges,
even though eight edges may have been obtained over the whole argument.

If $b>k/2$, the four edges satisfy the hypotheses of
Lemma~\ref{lem:twolarge}. Its three budget requirements follow from
\begin{align*}
 \ceil{k/2}+2m
   &\le\ceil{k/2}+2k-2T\le T,\\
 x+m&\le3k/2-T+1/2\le T,\\
 2x+b+\ceil{k/2}+2m
   &\le(9k+m-5T+4)/2\le3T.
\end{align*}
For the last inequality use $m\le k-T$ and $12T\ge10k+4$.
The first two use the second and third inequalities of
\eqref{eq:fourinteger}. The lemma produces a bad subfamily of at most seven
edges, completing the contradiction.
\end{proof}


\section*{Research assistance}
Claude (Anthropic) and Codex (OpenAI) were used for literature searches,
mathematical exploration, computational experiments, and manuscript
preparation. The proofs presented here are self-contained hand arguments;
no numerical output or software execution is required to verify them.

\begin{thebibliography}{9}
\bibitem{EFKT}
P.~Erd\H{o}s, D.~G.~Fon-Der-Flaass, A.~V.~Kostochka and Zs.~Tuza,
\emph{Small transversals in uniform hypergraphs},
Siberian Advances in Mathematics \textbf{2} (1992), 82--88.

\bibitem{FKW}
D.~G.~Fon-Der-Flaass, A.~V.~Kostochka and D.~R.~Woodall,
\emph{Transversals in uniform hypergraphs with property $(7,2)$},
Discrete Mathematics \textbf{207} (1999), 277--284.
\href{https://doi.org/10.1016/S0012-365X(99)00114-4}
{doi:10.1016/S0012-365X(99)00114-4}.
\bibitem{Kostochka}
A.~V.~Kostochka,
\emph{Transversals in uniform hypergraphs with property $(p,2)$},
Combinatorica \textbf{22} (2002), 275--285.
\href{https://doi.org/10.1007/s004930200013}{doi:10.1007/s004930200013}.

\bibitem{BKS}
M.~Buci\'c, D.~Kor\'andi and B.~Sudakov,
\emph{Covering graphs by monochromatic trees and Helly-type results
for hypergraphs}, Combinatorica \textbf{41} (2021), 319--352.
\href{https://doi.org/10.1007/s00493-020-4292-9}
{doi:10.1007/s00493-020-4292-9}.
\end{thebibliography}
\end{document}
